% GENERATED FILE. Edit canonical lessons, then run npm run generate:books.
% canonical-source-hash: fnv1a64:3751ab71a63d8f51
% canonical-lessons: TE-C474-sikharam, TE-C474-pithabhumi, TE-C474-khandam, TE-C474-agniparvatam, TE-C474-prasantamaina

\chapter{Mountains and Plateaus}
\label{ch:te-mountains-and-plateaus}
\hlchaptermodality{474}

\begin{chapteropening}
\textbf{By the end of this chapter:} \emph{I can say a peak, a plateau, a continent, a volcano and peaceful.}

Complete the last lesson of chapter 474: I can say a peak, a plateau, a continent, a volcano and peaceful.
\end{chapteropening}

\section[\texorpdfstring{\te{శిఖరం} (śikharaṁ)}{శిఖరం (śikharaṁ)}]{\te{శిఖరం} (śikharaṁ) — a peak; a summit}

\label{lesson:TE-C474-sikharam}



\begin{quote}
Before the new one: say the Telugu for a tidal surge, then the Telugu for to roar.
\end{quote}

\subsection*{You'll want to know: \te{శిఖరం}}
\textbf{\te{శిఖరం}} — \emph{śikharaṁ} — \textquotedblleft{}a peak; a summit\textquotedblright{}.

\begin{quote}
\te{ఉదయమే} \te{మేము} \te{కొండ} \te{శిఖరం} \te{చేరుకున్నాం}.
\end{quote}

\begin{quote}
\emph{udayamē mēmu koṇḍa śikharaṁ cērukunnāṁ.}
\end{quote}

\begin{quote}
Early in the morning we reached the hilltop.
\end{quote}

The top of a hill or a mountain, and the top of a temple tower.

\begin{cousinweb}
\textbf{\te{శిఖరం}} (\emph{śikharaṁ}) is borrowed from Sanskrit \emph{śikhara}, \textquotedblleft{}a point, a peak\textquotedblright{} (CDIAL 12435).
\end{cousinweb}

\begin{grammarlens}[title={What you've built}]
One more word for getting around.
\end{grammarlens}

\subsection*{Your turn}
\begin{itemize}
\raggedright
  \item \emph{Say it:} \emph{śikharaṁ}
  \item \emph{Say it:} \emph{śikharaṁ}, once more
  \item \emph{Say it:} \emph{udayamē mēmu koṇḍa śikharaṁ cērukunnāṁ.} — Early in the morning we reached the hilltop.
  \item \emph{Recall:} say the Telugu for a tidal surge, then the Telugu for to roar, then say \emph{śikharaṁ} again
\end{itemize}

\begin{tcolorbox}[breakable,colback=teal!4,colframe=teal!35!black,title={Before you move on}]
What does \te{శిఖరం} mean? (\textquotedblleft{}A peak; a summit\textquotedblright{}.) Say it once more.
\end{tcolorbox}
\section[\texorpdfstring{\te{పీఠభూమి} (pīṭhabhūmi)}{పీఠభూమి (pīṭhabhūmi)}]{\te{పీఠభూమి} (pīṭhabhūmi) — a plateau}

\label{lesson:TE-C474-pithabhumi}



\begin{quote}
Before the new one: say the Telugu for to roar, then the Telugu for a peak.
\end{quote}

\subsection*{You'll want to know: \te{పీఠభూమి}}
\textbf{\te{పీఠభూమి}} — \emph{pīṭhabhūmi} — \textquotedblleft{}a plateau\textquotedblright{}.

\begin{quote}
\te{ఆ} \te{పీఠభూమి} \te{మీద} \te{వర్షం} \te{తక్కువ}.
\end{quote}

\begin{quote}
\emph{ā pīṭhabhūmi mīda varṣaṁ takkuva.}
\end{quote}

\begin{quote}
There is little rain on that plateau.
\end{quote}

\te{పీఠ}, a seat, a stool, + \te{భూమి}, land: flat land raised up like a stool.

\begin{cousinweb}
\textbf{\te{పీఠభూమి}} (\emph{pīṭhabhūmi}) is a compound of two words borrowed from Sanskrit: \emph{pīṭha}, \textquotedblleft{}a stool, a bench\textquotedblright{} (CDIAL 8222), and \emph{bhūmi}, \textquotedblleft{}earth, ground\textquotedblright{} (CDIAL 9557).
\end{cousinweb}

\begin{grammarlens}[title={What you've built}]
One more word for getting around.
\end{grammarlens}

\subsection*{Your turn}
\begin{itemize}
\raggedright
  \item \emph{Say it:} \emph{pīṭhabhūmi}
  \item \emph{Say it:} \emph{pīṭhabhūmi}, once more
  \item \emph{Say it:} \emph{ā pīṭhabhūmi mīda varṣaṁ takkuva.} — There is little rain on that plateau.
  \item \emph{Recall:} say the Telugu for to roar, then the Telugu for a peak, then say \emph{pīṭhabhūmi} again
\end{itemize}

\begin{tcolorbox}[breakable,colback=teal!4,colframe=teal!35!black,title={Before you move on}]
What does \te{పీఠభూమి} mean? (\textquotedblleft{}A plateau\textquotedblright{}.) Say it once more.
\end{tcolorbox}
\section[\texorpdfstring{\te{ఖండం} (khaṇḍaṁ)}{ఖండం (khaṇḍaṁ)}]{\te{ఖండం} (khaṇḍaṁ) — a continent}

\label{lesson:TE-C474-khandam}



\begin{quote}
Before the new one: say the Telugu for a peak, then the Telugu for a plateau.
\end{quote}

\subsection*{You'll want to know: \te{ఖండం}}
\textbf{\te{ఖండం}} — \emph{khaṇḍaṁ} — \textquotedblleft{}a continent\textquotedblright{}.

\begin{quote}
\te{ప్రపంచంలో} \te{ఏడు} \te{ఖండాలు} \te{ఉన్నాయి}.
\end{quote}

\begin{quote}
\emph{prapañcaṁlō ēḍu khaṇḍālu unnāyi.}
\end{quote}

\begin{quote}
There are seven continents in the world.
\end{quote}

Plural: \te{ఖండాలు}, \emph{khaṇḍālu}.

Plural \te{ఖండాలు}. It also means a piece, a section.

\begin{cousinweb}
\textbf{\te{ఖండం}} (\emph{khaṇḍaṁ}) is borrowed from Sanskrit \emph{khaṇḍa}, \textquotedblleft{}broken; a piece, a part\textquotedblright{} (CDIAL 3792): a continent is one great piece of the earth.
\end{cousinweb}

\begin{grammarlens}[title={What you've built}]
One more word for getting around.
\end{grammarlens}

\subsection*{Your turn}
\begin{itemize}
\raggedright
  \item \emph{Say it:} \emph{khaṇḍaṁ}
  \item \emph{Say it:} \emph{khaṇḍaṁ}, once more
  \item \emph{Say it:} \emph{prapañcaṁlō ēḍu khaṇḍālu unnāyi.} — There are seven continents in the world.
  \item \emph{Recall:} say the Telugu for a peak, then the Telugu for a plateau, then say \emph{khaṇḍaṁ} again
\end{itemize}

\begin{tcolorbox}[breakable,colback=teal!4,colframe=teal!35!black,title={Before you move on}]
What does \te{ఖండం} mean? (\textquotedblleft{}A continent\textquotedblright{}.) Say it once more.
\end{tcolorbox}
\section[\texorpdfstring{\te{అగ్నిపర్వతం} (agniparvataṁ)}{అగ్నిపర్వతం (agniparvataṁ)}]{\te{అగ్నిపర్వతం} (agniparvataṁ) — a volcano}

\label{lesson:TE-C474-agniparvatam}



\begin{quote}
Before the new one: say the Telugu for a plateau, then the Telugu for a continent.
\end{quote}

\subsection*{You'll want to know: \te{అగ్నిపర్వతం}}
\textbf{\te{అగ్నిపర్వతం}} — \emph{agniparvataṁ} — \textquotedblleft{}a volcano\textquotedblright{}.

\begin{quote}
\te{ఆ} \te{ద్వీపంలో} \te{ఒక} \te{అగ్నిపర్వతం} \te{ఉంది}.
\end{quote}

\begin{quote}
\emph{ā dvīpaṁlō oka agniparvataṁ undi.}
\end{quote}

\begin{quote}
There is a volcano on that island.
\end{quote}

\te{అగ్ని}, fire, + \te{పర్వతం}, a mountain: a fire mountain.

\begin{cousinweb}
\textbf{\te{అగ్నిపర్వతం}} (\emph{agniparvataṁ}) is a compound of two words borrowed from Sanskrit: \emph{agni}, \textquotedblleft{}fire\textquotedblright{}, and \emph{parvata}, \textquotedblleft{}a rock, a mountain\textquotedblright{} (CDIAL 7945).
\end{cousinweb}

\begin{grammarlens}[title={What you've built}]
One more word for getting around.
\end{grammarlens}

\subsection*{Your turn}
\begin{itemize}
\raggedright
  \item \emph{Say it:} \emph{agniparvataṁ}
  \item \emph{Say it:} \emph{agniparvataṁ}, once more
  \item \emph{Say it:} \emph{ā dvīpaṁlō oka agniparvataṁ undi.} — There is a volcano on that island.
  \item \emph{Recall:} say the Telugu for a plateau, then the Telugu for a continent, then say \emph{agniparvataṁ} again
\end{itemize}

\begin{tcolorbox}[breakable,colback=teal!4,colframe=teal!35!black,title={Before you move on}]
What does \te{అగ్నిపర్వతం} mean? (\textquotedblleft{}A volcano\textquotedblright{}.) Say it once more.
\end{tcolorbox}
\section[\texorpdfstring{\te{ప్రశాంతమైన} (praśāntamaina)}{ప్రశాంతమైన (praśāntamaina)}]{\te{ప్రశాంతమైన} (praśāntamaina) — peaceful; calm}

\label{lesson:TE-C474-prasantamaina}



\begin{quote}
Before the new one: say the Telugu for a continent, then the Telugu for a volcano.
\end{quote}

\subsection*{You'll want to know: \te{ప్రశాంతమైన}}
\textbf{\te{ప్రశాంతమైన}} — \emph{praśāntamaina} — \textquotedblleft{}peaceful; calm\textquotedblright{}.

\begin{quote}
\te{కొండల} \te{మధ్య} \te{ఒక} \te{ప్రశాంతమైన} \te{గ్రామం} \te{ఉంది}.
\end{quote}

\begin{quote}
\emph{koṇḍala madhya oka praśāntamaina grāmaṁ undi.}
\end{quote}

\begin{quote}
There is a peaceful village among the hills.
\end{quote}

\te{ప్రశాంతంగా} is the adverb: \te{ప్రశాంతంగా} \te{నిద్రపోయాడు}, he slept peacefully.

\begin{cousinweb}
\textbf{\te{ప్రశాంతమైన}} (\emph{praśāntamaina}) is built from Sanskrit \emph{praśānta}, \textquotedblleft{}calmed, tranquil\textquotedblright{} (Monier-Williams), with the Telugu adjective ending -\te{మైన}.
\end{cousinweb}

\begin{grammarlens}[title={What you've built}]
One more word for getting around.
\end{grammarlens}

\subsection*{Your turn}
\begin{itemize}
\raggedright
  \item \emph{Say it:} \emph{praśāntamaina}
  \item \emph{Say it:} \emph{praśāntamaina}, once more
  \item \emph{Say it:} \emph{koṇḍala madhya oka praśāntamaina grāmaṁ undi.} — There is a peaceful village among the hills.
  \item \emph{Recall:} say the Telugu for a continent, then the Telugu for a volcano, then say \emph{praśāntamaina} again
\end{itemize}

\begin{tcolorbox}[breakable,colback=teal!4,colframe=teal!35!black,title={Before you move on}]
What does \te{ప్రశాంతమైన} mean? (\textquotedblleft{}Peaceful; calm\textquotedblright{}.) Say it once more.
\end{tcolorbox}
